\documentclass[10pt,a4paper]{amsart}
\usepackage[margin=19mm]{geometry}
\usepackage{amsmath,amssymb,mathtools}
\usepackage[scr=rsfs,cal=euler]{mathalfa}
\usepackage{xfrac}
\usepackage[unicode,hidelinks,bookmarks=false]{hyperref}
\newcommand{\ie}{i.e.\ }
\newcommand{\cf}{cf.\ }
\newcommand{\resp}{respectively\ }
\newcommand{\Subsection}{\subsection}
\providecommand{\nobreakdash}{\nobreak\hbox{-}}
\setlength{\emergencystretch}{2em}
\multlinegap0pt
\usepackage{bm}
\usepackage{bbm}
\usepackage{dsfont}
\usepackage{xspace}
\usepackage{cancel}
\usepackage{mathtools}
\usepackage{amsthm}
\makeatletter
\@ifundefined{theorem}{\newtheorem{theorem}{Theorem}}{}
\@ifundefined{lemma}{\newtheorem{lemma}[theorem]{Lemma}}{}
\@ifundefined{proposition}{\newtheorem{proposition}[theorem]{Proposition}}{}
\makeatother
\DeclareMathOperator{\first}{{first}}
\DeclareMathOperator{\last}{{last}}
\DeclareMathOperator{\val}{val}
\def\CN{{\mathcal N}}
\title[HNN extensions of free groups with equal associated subgroup]{HNN extensions of free groups with equal associated subgroups of finite index: polynomial time word problem}
\author{Hanwen Shen, Alexander Ushakov}
\date{Source: arXiv:2510.03801v2 (CC BY 4.0)}
\begin{document}
\maketitle

\noindent\emph{Source and licence.} HNN extensions of free groups with equal associated subgroups of finite index: polynomial time word problem. Version arXiv:2510.03801v2 (CC BY 4.0), distributed under the Creative Commons Attribution 4.0 International licence, as recorded in the arXiv licence metadata for this exact version and shown on the abstract page https://arxiv.org/abs/2510.03801v2. Full rights notice: licenses/arxiv-2510.03801-CC-BY-4.0.txt.\par\smallskip

\noindent\textit{Editorial scope.} Counted unit: the Proposition whose source label is \texttt{pr:main-step-case-I} in the article (source0.tex lines 1247--1269, source label \texttt{pr:main-step-case-I}), delivered together with the article's own complete original proof of it (source0.tex lines 1271--1335, one \texttt{proof} environment of 65 lines). No mathematics is written, repaired or re-derived: statement and proof are the article's own text line for line. Required context retained uncounted: le:first-last (lines 933--936); le:SLP-connected (lines 1031--1034); pr:Suv-SLP (lines 1059--1074); eq:slp-word (lines 1166--1170); pr:apply-phi (lines 1190--1209). Every definition, display and label of those lines is retained unchanged. Omitted from this package and not redistributed: 1--932, 937--1030, 1035--1058, 1075--1165, 1171--1189, 1210--1246, 1270--1270, 1336--1974. Nothing inside a retained excerpt is omitted or altered: every retained body is delivered exactly as the article's lines are written, so no omission receipt of any kind accompanies this package. Citations: the retained excerpts carry no citation command at all, so no bibliography entry of the article is transcribed and no reference is redistributed. Reference adaptation: none was needed and none was made; every reference inside the delivered unit resolves inside it, so no unresolved reference or citation is produced. Printed slips kept literal: no printed word, symbol, hypothesis or number of the counted statement or of its proof was altered. Layout and class: the article's own class is \texttt{amsart} with packages including url, fullpage, amssymb, latexsym, amsmath, mathrsfs, xy, enumerate, amsthm, psfrag, ifthen, comment, caption, todonotes; the wrapper is the lane's pinned \texttt{amsart} editorial wrapper at 10pt on A4, which restates the theorem-like environments the retained text needs and adds the article's own macro definitions that the retained text needs (\texttt{\textbackslash CN}, \texttt{\textbackslash first}, \texttt{\textbackslash last}, \texttt{\textbackslash val}), with extra wrapper packages (bm, bbm, dsfont, xspace, cancel, mathtools). The delivered numbering is the unit's own. Assets: no retained span contains a figure environment, a graphics-inclusion command, a PDF-page inclusion, a table or a TikZ picture; no image file is copied into this package, the package declares no asset dependency and no image file is redistributed with it. Licence: version arXiv:2510.03801v2 of this article is distributed under the Creative Commons Attribution 4.0 International licence, as recorded in the arXiv licence metadata for this exact version and shown on the abstract page https://arxiv.org/abs/2510.03801v2. A full rights notice is delivered at licenses/arxiv-2510.03801-CC-BY-4.0.txt. Characters: every retained excerpt is pure ASCII, so the delivered engine, XeLaTeX with its default Latin Modern text font, reports no missing character.
\begin{lemma}\label{le:first-last}
Given an SLP $P=(X,\CN,R,\delta)$, it takes $O(|P|)$ time to compute
$\first(R)$ and $\last(R)$.
\end{lemma}

\begin{lemma}\label{le:SLP-connected}
It takes linear time to decide if the sequence of edges $\val(P)$
is a path.
\end{lemma}

\begin{proposition}\label{pr:Suv-SLP}
Let $\Gamma=(V,E^\pm,\mu,r)$ be an $X$-regular 
and self-similar subgroup graph.
Let $u,v\in V$ and $S_{u,v}:E^\pm\to E^\pm$ be a 
permutation on the set of edges given explicitly
as a set of pairs $(e,S_{u,v}(e))$.
Given an SLP $P$ over $\Gamma$ 
it requires $O(|P|)$ time to construct an SLP $P'$
satisfying
\begin{itemize}
\item 
$\val(P')=S_{u,v}(\val(P))$,
\item 
$|P'|=|P|$.
\end{itemize}
\end{proposition}

\begin{equation}\label{eq:slp-word}
P_0, t^{\varepsilon_1},
P_1, \dots,
P_{k-1}, t^{\varepsilon_k},P_k
\end{equation}

\begin{proposition}\label{pr:apply-phi}
Let $P$ be an SLP over $\Gamma$ such that $\val(P)$ is a 
circuit in $\Gamma$ from $r$ to $r$.
Then $\mu(\val(P)) \in H$ and $\varphi$ is applicable to
$\mu(\val(P))$. 
There is an algorithm that
for a given $P$ produces an SLP $P'$ in $O(|P|)$ time
satisfying the following conditions:
\begin{itemize}
\item[(a)]
$\val(P')$ is a circuit based at $r$;
\item[(b)]
$\mu(\val(P')) = \varphi(\mu(\val(P)))$;
\item[(c)]
$|P'|\le |P| + \sum_{e\in E^+\setminus T}
|P_e^\varphi|+
|P_{e^{-1}}^\varphi|$.
\end{itemize}
The same holds for $\varphi^{-1}$.
\end{proposition}

\begin{proposition}\label{pr:main-step-case-I}
Consider a segment $P_{i-1},t^{-1},P_i,t,P_{i+1}$ in \eqref{eq:slp-word}.
It requires $O(|P_{i-1}|+|P_{i}|+|P_{i+1}|)$ time to check if
$P_i$ defines an element in $H$ 
(i.e., if $\mu(\val(P_i)) \in H$)
and, if so, to construct
an SLP $P$ satisfying the following properties:
\begin{itemize}
\item[(a)]
$\val(P)$ is a path in $\Gamma$ that starts at $r$,
\item[(b)]
$\mu(\val(P))\ =_G\ 
\mu(\val(P_{i-1})) 
\cdot
t^{-1} \mu(\val(P_{i})) t
\cdot
\mu(\val(P_{i+1}))$,
\item[(c)]
$|P| \le |P_{i-1}| + |P_i| + |P_{i+1}| + C$.
\end{itemize}
The same holds for segments $P_{i-1},t,P_i,t^{-1},P_{i+1}$, when
$P_i$ defines an element in $H$.
\end{proposition}

\begin{proof}
By Lemmas \ref{le:first-last} and \ref{le:SLP-connected},
we can check if $\val(P_i)$ defines a path that
starts and ends at $r$ (i.e., if $\mu(\val(P_i)) \in H$)
in $O(|P_i|)$ time.
By Proposition \ref{pr:apply-phi}, 
in $O(|P_i|)$ time we can compute an SLP $P_i'$ satisfying
\begin{itemize}
\item 
$\val(P_i')$ a circuit in $\Gamma$ from $r$ to $r$, and
\item 
$\mu(\val(P_i')) =_G t^{-1} \mu(\val(P_i)) t$.
\end{itemize}
Hence, the word
$$
\mu(\val(P_{i-1}))\mu(\val(P_{i}'))\mu(\val(P_{i+1}))
$$
defines the same element as the right-hand side of (b).
Note that, in general,
$\val(P_{i-1}) \circ \val(P_{i}') \circ \val(P_{i+1}))$
does not define a continuous path in $\Gamma$; 
there may be up to two points of discontinuity.

To create a required SLP $P$, it remains to 
properly concatenate the paths
$\val(P_{i-1})$, $\val(P_{i}')$, $\val(P_{i+1})$,
which can be done using shift operations.
Use Lemma \ref{le:first-last} and Proposition \ref{pr:Suv-SLP}
to compute 
$$
v = t(\last(P_{i-1})),
\ P_i''=S_{r,v}(P_i')
\ \mbox{ and }\ 
P_{i+1}'=S_{r,v}(P_{i+1}).
$$
Observe that $\val(P_i'')$ is a circuit and $t(\last(P_{i}'')) = v$.
Hence, by definition of $P_i''$ and $P_{i+1}'$,
the concatenation
$\val(P_{i-1}) \circ \val(P_i'') \circ \val(P_{i+1}')$
is a path in $\Gamma$ that starts at $r$.
Its label defines the same element as the right-hand side of (b)
because shift operations preserve labels.
Therefore, the SLP $P = P_{i-1} \circ P_i''\circ P_{i+1}'$
satisfies (a) and (b).

Finally, by construction,
$$
|P_i''| = |P_i'| \le |P_i| + \sum_{e\in E^+\setminus T}
|P_e^\varphi|+
|P_{e^{-1}}^\varphi|
\ \mbox{ and }\ 
|P_{i+1}'| = |P_{i+1}|
$$
because applying shift operations does not change the size of an SLP.
Thus, 
$$
|P| \le 
|P_{i-1}|+|P_{i}| + |P_{i+1}|+
\sum_{e\in E^+\setminus T}
|P_e^\varphi|+
|P_{e^{-1}}^\varphi| +2
\le |P_{i-1}| + |P_{i}| + |P_{i+1}| + C
$$
and $P$ satisfies (c).
\end{proof}

\end{document}
